\section{Preliminaries and notation}

This section establishes notation and standard definitions used throughout
the paper. No theorem specific to the graph

\[
X=\KC(L(P))
\]

is proved here.

\subsection{The Petersen graph and its line graph}

Let \(P\) denote the Petersen graph. Throughout this paper we regard \(P\)
as a fixed simple undirected graph with

\[
|V(P)|=10,
\qquad
|E(P)|=15,
\]

in which every vertex has degree three.

We write

\[
Y=L(P)
\]

for the line graph of \(P\). Thus the vertices of \(Y\) are the edges of
\(P\), and two vertices of \(Y\) are adjacent precisely when the
corresponding edges of \(P\) share a common endpoint.

\begin{proposition}
The line graph \(Y=L(P)\) is a connected quartic graph on fifteen vertices
and thirty edges.
\end{proposition}

\begin{proof}
Since the Petersen graph has fifteen edges,

\[
|V(Y)|=|E(P)|=15.
\]

Each edge of the cubic graph \(P\) is incident with two vertices, each
contributing two adjacent edges distinct from the edge itself.
Consequently every vertex of \(Y\) has degree four.

By the Handshake Lemma,

\[
|E(Y)|
=
\frac{15\cdot4}{2}
=
30.
\]

Because the Petersen graph is connected, its line graph is also connected.
\end{proof}

The natural action of \(\Aut(P)\) on the edge set induces an action on the
vertex set of \(Y\). Accordingly,

\[
\Aut(Y)\cong\Aut(P)\cong S_5.
\]

Throughout the paper we fix one labeled realization of the Petersen graph.
All explicit computational data refer to this fixed labeling. The complete
labeling is recorded in Appendix~A.

\subsection{Canonical bipartite double covers}

Let \(Y\) be a finite simple graph.

The canonical bipartite double cover of \(Y\), denoted
\(\KC(Y)\), has vertex set

\[
V(\KC(Y))
=
V(Y)\times\Ftwo,
\]

where

\[
(u,\varepsilon)
\sim
(v,\delta)
\]

if and only if

\[
u\sim_Y v
\]

and

\[
\delta=\varepsilon+1
\pmod2.
\]

Equivalently,

\[
\KC(Y)\cong Y\times K_2.
\]

The projection

\[
\pi:\KC(Y)\rightarrow Y
\]

is a regular two-fold covering projection with deck involution

\[
\tau(u,\varepsilon)
=
(u,\varepsilon+1).
\]

\begin{proposition}
If \(Y\) is connected, then \(\KC(Y)\) is bipartite.
Moreover, \(\KC(Y)\) is connected if and only if \(Y\) is not bipartite.
\end{proposition}

\begin{proof}
Every edge joins vertices having opposite second coordinate.
The connectedness criterion is the standard characterization of canonical
bipartite double covers.
\end{proof}

Since \(L(P)\) contains triangles, it is non-bipartite.
Consequently,

\[
X=\KC(L(P))
\]

is connected.

\subsection{Regular two-fold quotients}

Let \(Z\) be a connected graph. A \emph{pair partition} of \(V(Z)\) is a
partition of the vertex set into unordered pairs.

Every pair partition determines a unique fixed-point-free involution,
called its \emph{block swap}, which exchanges the two vertices in each
block. The block swap need not be an automorphism.

\begin{definition}
A \emph{regular two-fold quotient frame} of \(Z\) is a pair partition whose
block swap
\[
\tau\in\Aut(Z)
\]
is an involution and whose orbit projection
\[
\pi:
Z\longrightarrow Z/\langle\tau\rangle
\]
is locally bijective.
\end{definition}

Equivalently, the restriction of \(\pi\) to the neighborhood of every
vertex is a bijection onto the neighborhood of its image. In that case
\(\pi\) is a regular two-fold covering projection and
\(\langle\tau\rangle\) is its deck group.

Throughout this paper, \emph{quotient frame} means a regular two-fold
quotient frame in this sense.

\subsection{Group actions, orbitals, and coset graphs}

Let a finite group \(G\) act transitively on a finite set \(\Omega\).
The induced diagonal action on
\[
\Omega\times\Omega
\]
is given by
\[
(\alpha,\beta)^g
=
(\alpha^g,\beta^g).
\]

An \emph{orbital} is an orbit of this diagonal action. The orbitals
partition \(\Omega\times\Omega\) and form the basis relations of the
Schurian coherent configuration associated with the permutation action.
A permutation that permutes these basis relations is called a weak
automorphism of the configuration
\cite{Higman1975CoherentConfigurations,
Cameron2003CoherentConfigurations,
Eberhard2023HammingSandwiches}.

For an orbital \(R\), define its paired orbital by
\[
R^{*}
=
\{
(\beta,\alpha):
(\alpha,\beta)\in R
\}.
\]
The orbital \(R\) is \emph{self-paired} when \(R=R^{*}\). Every
self-paired non-diagonal orbital determines a simple undirected graph on
\(\Omega\), called its orbital graph.

Now let \(H\leq G\), let
\[
\Omega=G/H
\]
be the left coset space, and let \(G\) act by left multiplication. The
orbital containing \((H,gH)\) is determined by the double coset \(HgH\).
Its valency is
\[
[H:H\cap gHg^{-1}].
\]
It is self-paired precisely when
\[
HgH=Hg^{-1}H.
\]

Consequently, quartic orbital graphs correspond to self-paired double
cosets satisfying
\[
[H:H\cap gHg^{-1}]=4.
\]

\subsection{Centralizers of transitive coset actions}

Let \(G\) be a finite group, let \(H\leq G\), and let
\[
\Omega=G/H
\]
be the left coset space with the natural left multiplication action.

\begin{lemma}
\label{lem:coset-centralizer}
The centralizer of this transitive permutation representation is
\[
C_{\Sym(G/H)}(G)
\cong
N_G(H)/H.
\]
\end{lemma}

\begin{proof}
For \(n\in N_G(H)\), define
\[
\rho_n(xH)
=
xn^{-1}H.
\]
This is well-defined because \(n\) normalizes \(H\), and it commutes with
left multiplication by \(G\). The resulting homomorphism
\[
N_G(H)\longrightarrow C_{\Sym(G/H)}(G)
\]
has kernel \(H\).

Conversely, let \(c\) centralize the left action and write
\[
c(H)=xH.
\]
For every \(h\in H\),
\[
hxH
=
h\,c(H)
=
c(hH)
=
c(H)
=
xH.
\]
Thus \(x^{-1}Hx\leq H\). Since \(G\) is finite, equality holds, so
\(x\in N_G(H)\). Moreover,
\[
c(gH)
=
g\,c(H)
=
gxH
\]
for every \(g\in G\). Hence every element of the centralizer arises from
the construction above.
\end{proof}

The following standard observation records how normalizers interact with
orbital relations.

\begin{lemma}
\label{lem:normalizer-permutes-orbitals}
Let \(G\leq\Sym(\Omega)\), let
\[
n\in N_{\Sym(\Omega)}(G),
\]
and let \(R\) be an orbital of \(G\). Then
\[
R^n
=
\{
(\alpha^n,\beta^n):
(\alpha,\beta)\in R
\}
\]
is also an orbital of \(G\).

If \(R\) is self-paired, then \(R^n\) is self-paired, and \(n\) is an
isomorphism between the corresponding orbital graphs.
\end{lemma}

\begin{proof}
Choose \((\alpha,\beta)\in R\). Then
\[
R
=
\{
(\alpha^g,\beta^g):
g\in G
\}.
\]
Since \(n\) normalizes \(G\),
\[
R^n
=
\{
((\alpha^n)^{n^{-1}gn},
(\beta^n)^{n^{-1}gn}):
g\in G
\}
=
(\alpha^n,\beta^n)^G.
\]
Thus \(R^n\) is an orbital.

Coordinate reversal commutes with the diagonal action of \(n\), so
self-pairedness is preserved. The final assertion follows because \(n\)
relabels the vertices and carries \(R\) onto \(R^n\).
\end{proof}

\subsection{The two Klein four classes in \(S_5\)}

Up to conjugacy, \(S_5\) contains two classes of Klein four subgroups. The
first is
\[
V_{4,\mathrm{even}}
=
\{
1,
(12)(34),
(13)(24),
(14)(23)
\},
\]
whose nonidentity elements are all double transpositions. The second is
\[
V_{4,\mathrm{mixed}}
=
\langle(12),(34)\rangle
=
\{
1,
(12),
(34),
(12)(34)
\}.
\]

\begin{proposition}
The subgroups \(V_{4,\mathrm{even}}\) and
\(V_{4,\mathrm{mixed}}\) are not conjugate in \(S_5\).
\end{proposition}

\begin{proof}
Conjugation preserves cycle type. Every nonidentity element of
\(V_{4,\mathrm{even}}\) is a double transposition, whereas
\(V_{4,\mathrm{mixed}}\) contains transpositions.
\end{proof}

Both subgroups have index thirty, but their coset actions are not
permutation isomorphic.

\subsection{Binary voltage assignments and switching}

Let \(Y\) be a connected graph. A \emph{binary voltage assignment} is a
function
\[
\phi:E(Y)\longrightarrow\Ftwo.
\]
Its derived graph has vertex set
\[
V(Y)\times\Ftwo
\]
and lifted edges
\[
(u,\varepsilon)
\sim
(v,\varepsilon+\phi(uv))
\]
for every edge \(uv\in E(Y)\) and every \(\varepsilon\in\Ftwo\).

A switching function
\[
s:V(Y)\longrightarrow\Ftwo
\]
changes the assignment to
\[
\phi^{s}(uv)
=
\phi(uv)+s(u)+s(v).
\]
Switching-equivalent assignments define equivalent covers. Their switching
classes are naturally identified with
\[
H^1(Y,\Ftwo).
\]

For a cycle \(C\), its \emph{holonomy} is
\[
\phi(C)
=
\sum_{e\in E(C)}\phi(e)
\in\Ftwo.
\]
A cycle with holonomy \(0\) lifts to two cycles of the same length, while a
cycle with holonomy \(1\) lifts to one cycle of twice the length.

The constant-one assignment produces the canonical bipartite double cover.
Section~8 studies the \(S_5\)-fixed part of
\[
H^1(L(P),\Ftwo).
\]

\subsection{Notation}

Throughout the paper we use the following notation.

\begin{center}
\begin{tabular}{ll}
\toprule
Symbol & Meaning\\
\midrule
\(P\) & Petersen graph\\
\(Y=L(P)\) & Line graph of the Petersen graph\\
\(X=\KC(Y)\) & Canonical bipartite double cover of \(Y\)\\
\(A=\Aut(X)\) & Full automorphism group of \(X\)\\
\(K\cong S_5\) & Homogeneous-space action on \(X\)\\
\(C=C_A(K)\) & Centralizer of \(K\) in \(A\)\\
\(H=\Peven\) & Even Klein four subgroup of \(S_5\)\\
\(\Omega=S_5/H\) & Degree-thirty homogeneous space\\
\(\mathcal P_0\) & Natural quotient frame\\
\(\frorbit=A\cdot\mathcal P_0\) & Orbit of the natural frame\\
\bottomrule
\end{tabular}
\end{center}
